\section{Datasets and Experimental Setup}
\label{sec:experimental_setup}

\subsection{Datasets}
The main paper uses three public ultrasound datasets: TN5000 for thyroid nodules~\cite{Zhang_2025}, BUSI for breast lesions~\cite{Al_Dhabyani_2020}, and the Annotated Ultrasound Liver (AUL) dataset for liver lesions~\cite{aul_dataset}. TN5000 is the largest benchmark and is used for the most detailed ablation analysis. BUSI and AUL are used to test whether the architecture remains competitive under additional dataset-specific public ultrasound settings. They are not used as external thyroid-to-breast or thyroid-to-liver generalization tests.

\begin{table}[t]
    \centering
    \small
    \caption{Overview of the main datasets used in the paper. BUSI and AUL are used in a benign--malignant lesion-classification setting.}
    \label{tab:datasets}
    \begin{adjustbox}{max width=\columnwidth}
    \begin{tabular}{llllll}
        \toprule
        Dataset & Organ & Images & B / M & Split protocol & Role \\
        \midrule
        TN5000 & Thyroid & 5000 & 1426 / 3574 & Official 3500/500/1000 train/val/test & Main \\
        BUSI & Breast & 647 & 437 / 210 & Fixed test 129; five training folds on trainval 518 & Additional \\
        AUL & Liver & 635 & 200 / 435 & Fixed test 127; five training folds on trainval 508 & Additional \\
        \bottomrule
    \end{tabular}
    \end{adjustbox}
\end{table}

\subsection{Implementation details}
All datasets are processed in an ROI-based manner, but the crop expansion differs across benchmarks. TN5000 and BUSI use a lesion expansion ratio of 0.30, whereas AUL uses a lighter expansion ratio of 0.20. The lesion bounding box is read from XML or manifest metadata, expanded with a minimum crop size of 64 pixels, clipped to the image boundary, and resized to $256 \times 256$; when a reliable crop is unavailable, the whole image is used as a fallback. The current protocol preserves the expanded box aspect ratio before resizing and does not force a square crop. The only input-size exception in the comparison pool is EfficientFormer-L1, which is evaluated at $224 \times 224$ because its wrapper is not compatible with the $256 \times 256$ ROI setting used by the other models.

The training pipeline uses the same recipe family across the main runs: random horizontal flip, random rotation ($\pm 4^\circ$), mild affine translation and scaling, brightness/contrast jitter, and random Gaussian blur. These perturbations are kept mild because ultrasound gray-scale intensity, speckle, and gain variation are part of the imaging domain rather than arbitrary color content. Images are loaded as RGB tensors for compatibility with ImageNet-initialized backbones; gray-scale inputs are converted to three channels by the image loader. Evaluation uses deterministic resizing and ImageNet normalization only. All models use AdamW with weight decay $10^{-4}$, a lightweight MLP head with dropout 0.3, and label smoothing 0.0. The learning-rate schedule is dataset-specific.

TN5000 uses the official train/val/test split and is reported with three random seeds ($17$, $27$, and $37$). The mainline setting uses 70 epochs, patience 12, manual class weights, and differential learning rates of $5 \times 10^{-5}$ for the backbone and $1.5 \times 10^{-4}$ for the classification head. BUSI uses a fixed held-out test set together with five training folds on the trainval subset; its main runs use 200 epochs, patience 36, class weighting, and the same $5 \times 10^{-5} / 1.5 \times 10^{-4}$ learning-rate pair. AUL likewise uses a fixed held-out test set plus five training folds on trainval, but uses a lighter budget: 20 epochs, patience 4, no class weighting, top-2 checkpoint ensembling, and lower differential learning rates of $1 \times 10^{-5}$ for the backbone and $3 \times 10^{-5}$ for the classification head. Thus, BUSI/AUL mean $\pm$ standard deviation values quantify variation across training folds evaluated on a fixed test set, not an outer five-fold cross-validation estimate where each case appears in a test fold exactly once. The mainline scripts enable exponential moving average (EMA) with decay 0.9995; the exception in the baseline pool is BUSI RepViT-M1.1, for which EMA is disabled because the EMA variant produced unstable folds during verification.

For inference-time calibration, each run uses validation-derived temperature scaling~\cite{guo2017calibration}, averages the top-$k$ validation-improved checkpoints (top-3 on TN5000 and BUSI, top-2 on AUL), applies horizontal-flip test-time augmentation, fits a scalar temperature on validation logits, and searches a decision threshold from 0.10 to 0.95 with step 0.01 using balanced accuracy as the default selection criterion. The resulting validation threshold is applied to the test predictions of that run. Main benchmark tables report mean $\pm$ standard deviation across the three TN5000 seeds or the five BUSI/AUL folds.

Balanced accuracy is used for threshold selection because it gives equal weight to benign and malignant recall and is less tied to class prevalence than raw accuracy. AUC is kept as the threshold-free ranking metric, while macro F1 and accuracy are reported as complementary operating-point summaries.

Importantly, all temperatures and decision thresholds are fixed strictly from validation-derived predictions. No test labels are used for temperature fitting, threshold tuning, or operating-point selection; the test set is used only once for final reporting after these choices are fixed.

\subsection{Automatic ROI closed-loop protocol}
The main experiments use annotation-derived ROI crops to isolate classifier behavior under a controlled input protocol. To evaluate whether this ROI classifier can be connected to an automatic front end at inference time, we additionally train one-class lesion detectors on BUSI and AUL and perform closed-loop classification with detector-predicted crops. The detector is used only as an auxiliary ROI generator; it is not presented as a novel detection method.

For each BUSI/AUL fold, a one-class YOLO11n detector is trained on the corresponding training annotations and evaluated on the held-out test images. A parallel TN5000 automatic-ROI check trains YOLO11n detectors under three detector seeds and evaluates the official test split. The predicted lesion box with the highest confidence is exported into a VOC-style annotation file and passed through the same expansion, resize, normalization, temperature-scaling, and thresholding code used by the ROI classifier. We then compare three test-time inputs using the same trained UL-TransXNet checkpoints: \textit{oracle ROI}, which uses the ground-truth lesion annotation; \textit{automatic ROI}, which uses the detector-predicted box; and \textit{full image}, which disables the lesion crop and resizes the complete ultrasound image. Thus, the closed-loop experiment tests whether detector-predicted crops can partially substitute annotation-derived ROI construction without retraining the classifier.

\subsection{Compared methods}
The benchmark pool combines representative CNN, lightweight hybrid, and modern efficient baselines. The comparison set includes ResNet50~\cite{he2016resnet}, EfficientNet-B0~\cite{tan2019efficientnet}, MobileNetV3-Large~\cite{howard2019mobilenetv3}, Swin-T~\cite{liu2021swin}, DenseNet121~\cite{huang2017densenet}, ConvNeXt-Tiny~\cite{liu2022convnext}, RepViT-M1.1~\cite{repvit}, and EfficientFormer-L1~\cite{li2022efficientformer}. The main paper table reports \textit{Ours} together with the strongest non-\textit{Ours} baseline on each dataset, where ``strongest'' is determined by mean test AUC under the same reporting convention. The full baseline table is provided in Appendix Table~\ref{tab:full_benchmark}.

\subsection{Baseline fairness and reporting convention}
All models are evaluated through the same ROI construction, label mapping, train/validation/test splits, metric computation, validation-derived threshold search, temperature-scaling procedure, and test-time augmentation convention unless an exception is explicitly stated. All backbone models use ImageNet-pretrained initialization when available. The classification head is replaced by the binary head used in this study, so ImageNet classifier logits are not used as medical features. The main exceptions are practical and are reported rather than hidden: EfficientFormer-L1 uses $224 \times 224$ inputs because its wrapper is not compatible with the $256 \times 256$ ROI setting used by the other models; BUSI RepViT-M1.1 disables EMA because the EMA variant produced unstable folds during verification; and AUL uses a lighter training budget than TN5000/BUSI because of the smaller and less stable training setting. These exceptions mean that the benchmark should be interpreted as a carefully matched protocol comparison, not as an exhaustive hyperparameter search for every baseline.

The main results table intentionally reports the proposed model and the strongest non-\textit{Ours} baseline per dataset for readability, while Appendix Table~\ref{tab:full_benchmark} gives the complete baseline pool. This prevents the main text from implying that a single baseline was chosen in advance as the only comparator. It also makes clear that ConvNeXt-Tiny remains the strongest TN5000 operating-point baseline, whereas UL-TransXNet is selected for its validation-audited three-dataset behavior.

\subsection{Evaluation metrics and statistical analysis}
The main metrics are AUC, balanced accuracy (BalAcc), macro F1, and accuracy. Unless explicitly stated otherwise, ``F1'' in the manuscript refers to macro F1. Because the tables emphasize repeated-seed/fold stability, results are reported as mean $\pm$ standard deviation. Thresholded metrics should be interpreted as validation-selected operating points rather than test-tuned oracle points.

As a complementary diagnostic analysis, we also construct case-level summaries from the frozen test prediction CSV files. For TN5000, the three seed predictions are averaged per test case; for BUSI and AUL, the five fixed-test fold predictions are averaged per test case. AUC, sensitivity, specificity, PPV, NPV, ECE, and Brier score are then computed on these averaged predictions, with 95\% confidence intervals estimated by bootstrap resampling over test cases. This analysis is used to expose operating-point and calibration behavior on the fixed test sets; it is not treated as prospective clinical significance evidence.

\subsection{Local efficiency profiling}
Model efficiency is profiled on the TN5000 test split, which provides the largest held-out set among the main datasets. Local CUDA latency is measured on an NVIDIA GeForce RTX 4080 SUPER using batch size 1. Each model is instantiated with the same binary-classification wrapper used for complexity profiling, warmed up for 50 forward passes, and then timed for 300 model-only forward passes with CUDA synchronization. We report the mean model-only latency and the corresponding FPS. A separate TN5000 pipeline pass over 1000 test ROI crops is also logged for diagnosis, but the main trade-off table uses model-only latency to avoid mixing architecture cost with image decoding and CPU preprocessing overhead.

\subsection{Mobile prototype protocol}
Mobile inference is evaluated with an Android prototype built on ONNX Runtime 1.24.3. The reported batch experiments were executed on a Xiaomi device (model 24129PN74C, Android 16, arm64-v8a). The prototype study uses TN5000-calibrated UL-TransXNet ONNX exports and compares three inference variants: a single checkpoint, the same checkpoint with horizontal-flip test-time augmentation (TTA), and a three-checkpoint ensemble with TTA. Cold-start runs explicitly cleared the session cache before evaluation, whereas warm-start runs reused the loaded model session. The single-checkpoint mode was measured with one cold-start and three warm-start runs; the TTA and ensemble variants were measured with one cold-start and one warm-start run each. The prototype configuration used lesion expansion ratio 0.30, a fitted temperature of 1.819830, and decision threshold 0.73 inherited from TN5000 validation calibration. We report average end-to-end latency and the corresponding 95th percentile latency over the full ROI preprocessing-plus-inference pipeline.
